\documentclass[12pt,a4paper]{article}
\usepackage{geometry}
\geometry{margin=2.5cm}
\usepackage{fontspec}
\usepackage{xepersian}
\settextfont{Arial}
\begin{document}

\section*{تطابق مدل ILP + گراف انقباض با فرمالیسم مدل‌های ترکیبیاتی}

فرمالیسم ۴-تایی (که معمولاً شامل فضای نمونه‌ها، فضای جواب‌های موجه، تابع هدف و جهت بهینه‌سازی است) یک تعریف انتزاعی و تئوریک (Theoretical) از مسائل بهینه‌سازی ترکیبیاتی است. اما ترکیب یک گراف انقباض (Contraction Graph) با یک مدل برنامه‌ریزی خطی عدد صحیح (ILP)، در واقع نمایش جبری و عملیاتی (Algebraic/Operational) دقیقاً همان فرمالیسم است.

با توجه به منطق مقاله‌ی (Doron-Arad et al. 2020) و مفاهیم پایه‌ای بهینه‌سازی، این دو به دلایل زیر کاملاً بر هم منطبق هستند و مدل ILP پیاده‌سازی شده در این پروژه یک مدل ترکیبیاتی کامل محسوب می‌شود:

\subsection*{۱. گراف انقباض = بازتعریف هوشمندانه‌ی فضای نمونه‌ها (Instance Space)}
در مسائل ترکیبیاتی پیچیده (مثل مقاله مذکور که با تضادهای گروهی Clique-graph conflicts سروکار دارد، و یا این پروژه که دارای محدودیت‌های Cohesive یا Must-link است)، گنجاندن تمام محدودیت‌ها در فرمول‌های ریاضی باعث انفجار فضای حالت می‌شود. استفاده از گراف انقباض (پیدا کردن مولفه‌های همبند و تبدیل آن‌ها به یک ابرگره یا Super-node)، در واقع یک پیش‌پردازش ساختاری است که فضای مسئله را به یک فضای معادل اما کوچک‌تر و خوش‌رفتارتر نگاشت می‌کند. در اینجا بخشی از محدودیت‌های قطعی در ساختار داده حل شده‌اند تا مدل نهایی تمیزتر شود.

\subsection*{۲. محدودیت‌های ILP = تعریف صریح فضای موجه (Feasible Space / Polytope)}
در فرمالیسم ۴-تایی، به یک مجموعه از جواب‌های موجه ($F$) نیاز است. در ILP، این فضای موجه به کمک معادلات و نامعادلات خطی ساخته می‌شود (مثل اینکه هر مولفه منقبض‌شده فقط در یک سطل قرار گیرد و ظرفیت سطل نقض نشود).
اگر به مقاله (Theorem 3.11) رجوع کنیم، نویسندگان برای اثبات وجود یک جواب موجه برای آیتم‌های کوچک، دقیقاً یک فضای موجه یا پلی‌توپ (Polytope) را با استفاده از محدودیت‌های ریاضی و متغیرهای تصمیم تعریف می‌کنند. محدودیت‌های ILP دقیقاً همین کار را می‌کنند: آن‌ها مرزهای هندسی فضای شدنی (Feasible Space) را به زبان جبر خطی می‌سازند.

\subsection*{۳. تابع هدف ILP = تابع هدف در فرمالیسم (Objective Function)}
متغیرهایی مثل $y$ در ILP و کمینه‌سازی جمع آن‌ها، دقیقاً همان تابع نگاشتی است که به هر نقطه از فضای موجه، یک هزینه (Cost) اختصاص می‌دهد و سالور به دنبال مینیمم کردن آن در مرزهای همان Polytope است.

\subsection*{نتیجه‌گیری}
فرمالیسم ۴-تایی، زبان نظریه مجموعه‌ها برای تعریف مسئله است، در حالی که ILP به همراه گراف انقباض، زبان جبر خطی و گراف برای مدل‌سازی دقیقاً همان مسئله است. گراف انقباض متولی مدیریت قیود ساختاری (Structural Constraints) است تا مدل کوچک شود، و متغیرهای باینری و محدودیت‌های ILP دقیقاً مرزهای فضای شدنی (Feasible Space) را می‌سازند و تابع هدف روی آن‌ها بهینه‌سازی می‌شود (مشابه رویکرد تعریف Polytope در مقالاتی مثل Doron-Arad et al). در نتیجه، این یک نمایش کاملاً استاندارد و فرمال از یک مدل بهینه‌سازی ترکیبیاتی است.

\end{document}
